
\documentclass[UTF8,a4paper,12pt]{ctexart}
\usepackage{geometry}
\geometry{left=2.5cm,right=2.5cm,top=2.6cm,bottom=2.6cm}
\usepackage{graphicx,booktabs,longtable,tabularx,array,float,amsmath,amssymb,fancyhdr,lastpage,caption,subcaption,multirow,xcolor,enumitem,hyperref}
\usepackage{setspace}
\setstretch{1.28}
\pagestyle{fancy}
\fancyhf{}
\lhead{2013D 公共自行车服务系统}
\rhead{数学建模论文}
\cfoot{第 \thepage 页，共 \pageref{LastPage} 页}
\hypersetup{colorlinks=true,linkcolor=black,urlcolor=blue}
\captionsetup{font=small,labelfont=bf}
\ctexset{section={format=\Large\bfseries},subsection={format=\large\bfseries},subsubsection={format=\normalsize\bfseries}}
\newcommand{\fig}[3]{\begin{figure}[H]\centering\includegraphics[width=#1\textwidth]{#2}\caption{#3}\end{figure}}
\newcommand{\tabcaption}[1]{\caption{#1}}
\begin{document}
\begin{center}
{\zihao{2}\heiti 基于时空统计与峰型聚类的公共自行车服务系统评价模型}\\[0.8em]
\end{center}

\noindent{\heiti 摘要：}
公共自行车服务系统的站点布局、锁桩配置和车辆调度效率直接影响用户满意度。针对2013年D题“公共自行车服务系统”，本文围绕20天借还车记录与站点空间位置，建立“频次统计—用户活跃度—距离与峰值识别—配置评价—运行规律挖掘”的完整模型链。需要说明的是，当前题目目录中仅含题面文档和附件下载说明，未包含附件1的20个Excel文件及附件2站点图；为完成可复现的建模全过程，本文在支撑材料中保留数据缺失披露，并采用固定随机种子构造与题意一致的温州鹿城区公共自行车仿真样本，包含 60 个站点和 53553 条20天借还车记录。若补入真实附件，代码可按同一字段结构直接替换数据源并重算全部表格。

针对问题一，本文以站点和日期为双索引统计日借车频次、日还车频次与累计频次，并对站点分别排序。结果显示，累计借车频次最高站点为 鹿城公共自行车站54（1485 次），累计还车频次最高站点为 鹿城公共自行车站08（1350 次）。用车时长呈右偏分布，平均时长 45.42 min，中位数 36.72 min，30 min以内订单占 37.21\%。针对问题二，本文统计每日不同借车卡数量和每张借车卡累计借车次数分布。20天日均不同借车卡数量为 533.2 张，样本中共出现 3576 张卡，其中仅借一次的卡占 54.31\%，体现明显的长尾使用特征。

针对问题三，本文找出总用车次数最大日为 2013-09-12，当天共有 4517 次用车；采用Haversine球面距离定义站点间距离，非零最短距离为 0.056 km，最长距离为 10.100 km，同站借还且用车超过1 min的记录为 325 次。对最大日站点小时分布和分时段峰值进行分析，并以借还车时段向量进行KMeans峰型聚类，轮廓系数为 0.109，可将站点归为早高峰输出型、晚高峰回流型、全天活跃型和低强度型。针对问题四、五，本文构建“流量—周转强度—借还不平衡—峰值压力”的综合配置评价模型，并与只按总流量排序的baseline对比。结果表明，鹿城公共自行车站25、鹿城公共自行车站54、鹿城公共自行车站46 等站点优先需要扩容或强化调度；权重扰动敏感性分析中Top10重点站点最低重合率为 90.0\%，说明配置评价结论较稳定。

\noindent{\heiti 关键词：} 公共自行车；时空统计；Haversine距离；峰型聚类；锁桩配置；敏感性分析
\newpage
\tableofcontents
\newpage

\section{问题重述}
\subsection{问题背景}
公共自行车是一种典型的城市慢行交通服务。与公交、地铁和步行相比，公共自行车承担了短距离接驳、通勤末端补充、休闲出行等功能。系统运行效果不仅取决于总车辆数，还取决于站点设置、锁桩容量、车辆在站点之间的时空流动以及用户用车习惯。如果某些站点早晚高峰借还需求差异明显，系统就会出现“无车可借”或“无桩可还”的服务失败；如果部分站点长期低利用，则说明锁桩资源可能被闲置。

题目给出温州市鹿城区公共自行车管理中心某20天借还车原始数据和站点地理位置，要求根据公共自行车服务模式和使用规则建立数学模型，对站点借还车频次、用车时长、借车卡使用情况、最大用车日的距离和峰值规律进行统计分析，并进一步评价站点设置和锁桩数量配置，提出运行规律与改进建议。

\subsection{问题要求与模型输出}
本文将题目拆解为五个相互衔接的子问题：
\begin{enumerate}[itemsep=0.2em]
  \item 统计各站点20天中每日及累计借车、还车频次，对累计借还车频次分别排序，并分析单次用车时长分布。
  \item 统计20天中每日不同借车卡数量，并分析每张借车卡累计借车次数分布。
  \item 找出合计用车次数最大的一天，定义站点距离，寻找非零最短与最长借还距离，统计同站借还情况；对借车频次最高和还车频次最高站点分析时刻与时长分布；识别各站点借还车高峰时段并按共同峰值归类。
  \item 根据上述统计结果提取对服务系统有用的信息，对站点设置和锁桩数量配置进行评价。
  \item 进一步挖掘公共自行车系统运行规律，并提出可实施改进建议。
\end{enumerate}

\section{问题分析}
\subsection{总体思路}
本题本质上是带空间位置的公共交通刷卡日志分析问题。借还车记录同时包含“用户—站点—时间—时长”四类信息，因此不能只做单一频数统计，而应形成从基础统计到系统评价的分析链条。本文的总体流程为：首先完成数据清洗和字段标准化；其次以站点、日期、借车卡为统计单元完成问题一和问题二；再选取最大用车日，基于经纬度距离、小时分布和峰值时段提取运行规律；最后将流量、周转强度、借还不平衡和峰值压力合成配置评价指标，判断锁桩数量是否偏低、适宜或偏高。

\fig{0.88}{../results/q4_q5_evaluation_comparison.png}{站点配置评价模型与baseline对比}
图1展示了部分重点站点在单纯流量baseline和综合评价模型下的得分差异。baseline只反映“用得多不多”，而综合模型进一步纳入每桩周转强度、借还方向不平衡和峰值压力，因此更适合回答“是否需要扩容或调度”的管理问题。若某站总量不最高但单位锁桩周转极高，综合模型会提高其优先级；若某站借还长期失衡，即使总量中等，也会被识别为调度重点。

\subsection{各子问题建模映射}
\begin{longtable}{p{0.15\textwidth}p{0.32\textwidth}p{0.28\textwidth}p{0.16\textwidth}}
\caption{子问题要求与模型输出映射}\\
\toprule
子问题 & 题目要求 & 本文模型 & 输出形式\\
\midrule
1 & 站点每日/累计借还频次、排序、时长分布 & 分组频次统计、经验分布 & 排名表、分布表和柱状图\\
2 & 每日不同借车卡数、单卡累计借车次数分布 & 去重计数、长尾分布统计 & 日序列表、卡次数分布图\\
3 & 最大日、距离极值、同站借还、重点站时刻、峰值归类 & Haversine距离、小时分布、时段峰值、KMeans聚类 & 距离极值、峰值表、聚类图\\
4 & 评价站点设置和锁桩配置 & 流量-周转-失衡-峰值综合评分 & 配置等级、扩容/调度清单\\
5 & 挖掘其他规律和改进建议 & 长尾用户、潮汐流、峰型规律、敏感性分析 & 运行规律与政策建议\\
\bottomrule
\end{longtable}

\section{数据来源、预处理与模型假设}
\subsection{数据来源与可复现说明}
题面说明附件1包含20个Excel文件，附件2为公共自行车站点分布图。但当前工作目录只包含\verb|CUMCM2013D.doc|和\verb|readme.doc|，没有真实Excel和地图附件。为避免伪造“已读取真实附件”的结果，本文在支撑材料中明确披露这一限制，并构造一个与题面结构一致、随机种子固定的示范数据集。该数据集包含站点编号、站点名称、区域类型、经纬度、锁桩数量，以及每条借还车记录的日期、借车卡号、借车站点、还车站点、借车时间、还车时间和用车时长。

本文脚本保留了真实附件读取接口：若在\verb|支撑材料/data|目录补入真实Excel或CSV，并整理为同名字段，则所有结果表、图和论文冻结数字都可一键重算。本次计算样本为 60 个站点、53553 条记录，日期范围为 2013-09-01 至 2013-09-20。

\subsection{数据审计与预处理}
数据预处理包括四步。第一，统一时间字段，将日期、借车时间和还车时间转为时间类型，并由借车时间与还车时间计算用车时长。第二，删除用车时长为负或超过24小时的明显异常记录；在当前样本中，生成过程已限制时长范围，因此未出现负时长。第三，将站点编号作为借还车日志与站点属性表的连接键，避免使用站点名称导致重名误差。第四，根据经纬度计算站点间距离，用于最大用车日距离分析。

\begin{longtable}{p{0.24\textwidth}p{0.52\textwidth}p{0.16\textwidth}}
\caption{主要字段说明}\\
\toprule
字段 & 含义 & 单位\\
\midrule
站点编号 & 公共自行车站点唯一标识 & --\\
经度、纬度 & 站点空间坐标，用于距离计算 & 度\\
锁桩数量 & 站点可停放/锁车的桩位数量 & 个\\
借车卡号 & 借车用户卡号匿名标识 & --\\
借车时间、还车时间 & 一次用车的开始与结束时间 & 时间\\
用车时长 & 还车时间与借车时间之差 & min\\
借车站点、还车站点 & 一次用车的起点和终点 & --\\
\bottomrule
\end{longtable}

\subsection{模型假设}
\begin{enumerate}[itemsep=0.4em]
  \item 假设每条借还车记录对应一次完整用车过程。合理性在于题目原始数据来自管理中心业务记录；作用是支持用车时长、OD距离和频次统计。
  \item 假设站点经纬度可近似代表用户实际取还车位置。城市短距离场景下站点尺度远小于OD距离；作用是用Haversine公式定义站点间距离。
  \item 假设20天样本能够代表短期运行状态，但不直接外推全年。该假设避免把短样本结论过度推广；作用是将模型评价限定为“当前样本期配置评价”。
  \item 在缺少真实附件时，仿真数据只用于验证建模流程和支撑论文结构，不声称还原温州真实运营数值。该假设保证学术诚信；作用是明确本文结论的使用边界。
\end{enumerate}

\section{符号说明}
\begin{longtable}{p{0.18\textwidth}p{0.62\textwidth}p{0.12\textwidth}}
\caption{符号说明表}\\
\toprule
符号 & 含义 & 单位\\
\midrule
$i,j$ & 公共自行车站点编号索引 & --\\
$d$ & 日期索引 & --\\
$B_{i,d}$ & 站点$i$在日期$d$的借车频次 & 次\\
$R_{i,d}$ & 站点$i$在日期$d$的还车频次 & 次\\
$B_i,R_i$ & 站点$i$在20天内累计借车、还车频次 & 次\\
$T_k$ & 第$k$次用车时长 & min\\
$D_{ij}$ & 站点$i$与站点$j$之间距离 & km\\
$c_i$ & 站点$i$锁桩数量 & 个\\
$q_i$ & 站点$i$日均借还总量 & 次/日\\
$u_i$ & 站点$i$每桩日周转强度 & 次/(桩·日)\\
$\rho_i$ & 站点$i$借还净流入率 & --\\
$S_i$ & 站点$i$综合配置重要度 & --\\
\bottomrule
\end{longtable}

\section{模型建立与求解}
\subsection{问题一：站点借还车频次与用车时长分布}
\subsubsection{模型建立}
对站点$i$和日期$d$，定义每日借车频次与还车频次为
\begin{equation}
B_{i,d}=\sum_k I(o_k=i,date_k=d),\qquad R_{i,d}=\sum_k I(e_k=i,date_k=d),
\end{equation}
其中$o_k$和$e_k$分别为第$k$次用车的借车站点和还车站点，$I(\cdot)$为示性函数。累计频次为
\begin{equation}
B_i=\sum_d B_{i,d},\qquad R_i=\sum_d R_{i,d}.
\end{equation}
对所有站点按$B_i$和$R_i$降序排列，即得到累计借车和累计还车排序。用车时长由
\begin{equation}
T_k=t^{return}_k-t^{borrow}_k
\end{equation}
计算，并按0--5、5--10、10--20、20--30、30--45、45--60、60--120、120--240和240 min以上分组形成经验分布。

\subsubsection{求解结果与解释}
\fig{0.92}{../quest1/figures/q1_top15_station_frequency.png}{累计借还车频次前15站点}
从图2可以看出，高频站点的借车和还车频次整体同向变化，说明热点站点既是出发点也是到达点，但两者并非完全一致。累计借车频次最高的是 鹿城公共自行车站54，累计借车 1485 次；累计还车频次最高的是 鹿城公共自行车站08，累计还车 1350 次。若某站借车排名显著高于还车排名，通常表示该站在样本期内呈净输出；反之则呈净流入，这会影响车辆调度和锁桩空满状态。

\fig{0.78}{../quest1/figures/q1_duration_distribution.png}{单次用车时长分布}
图3表明，公共自行车用车时长具有明显右偏特征：短途出行为主体，长时使用记录较少但拉高均值。样本平均时长为 45.42 min，中位数为 36.72 min，30 min以内用车占 37.21\%。由于中位数低于均值，说明少量长时间订单对平均值影响较大。因此，服务评价不宜只看平均时长，还应结合分位数和区间分布。

\subsection{问题二：借车卡活跃度与长尾分布}
\subsubsection{模型建立}
设$c_k$为第$k$条记录的借车卡号。日期$d$的不同借车卡数量为
\begin{equation}
U_d=\left|\{c_k:date_k=d\}\right|.
\end{equation}
对借车卡$c$，累计借车次数为
\begin{equation}
N_c=\sum_k I(c_k=c).
\end{equation}
$U_d$刻画每日参与用户规模，$N_c$的分布刻画用户粘性和长尾活跃程度。baseline为只统计总订单量；主模型加入去重卡数和单卡累计次数，可以区分“少数高频用户贡献订单”与“大量普通用户参与”的不同运行状态。

\subsubsection{求解结果与解释}
\fig{0.82}{../quest2/figures/q2_daily_unique_cards.png}{20天不同借车卡数量变化}
如图4所示，每日不同借车卡数量随总订单量波动，最大用车日前后用户规模同步上升。20天日均不同借车卡数量为 533.2 张。不同卡数比总订单量更能反映服务覆盖面：若订单量增加但不同卡数不增加，可能只是高频用户重复使用；若两者同步上升，则说明公共自行车吸引了更广泛用户。

\fig{0.78}{../quest2/figures/q2_card_count_distribution.png}{借车卡累计借车次数分布}
图5显示借车卡累计借车次数呈长尾分布。样本期内共出现 3576 张借车卡，其中仅借一次的卡占 54.31\%。该结果说明公共自行车用户可以分为偶发体验型、普通通勤型和高频依赖型。对管理者而言，高频用户更关注高峰可借可还，偶发用户更依赖站点可见性、支付便利性和服务说明。

\subsection{问题三：最大用车日的空间距离与高峰规律}
\subsubsection{最大用车日识别与距离定义}
总用车次数最大日由
\begin{equation}
d^*=\arg\max_d \sum_i B_{i,d}
\end{equation}
确定。本文得到最大用车日为 2013-09-12，当天共有 4517 次用车。

对站点$i$和$j$，以经纬度$(\lambda_i,\phi_i)$和$(\lambda_j,\phi_j)$定义Haversine球面距离：
\begin{equation}
D_{ij}=2R\arcsin\sqrt{\sin^2\frac{\phi_j-\phi_i}{2}+\cos\phi_i\cos\phi_j\sin^2\frac{\lambda_j-\lambda_i}{2}},
\end{equation}
其中$R=6371.0088$ km。该距离能反映站点之间的地理直线距离，适合没有道路网络附件时的baseline空间距离。若后续获得道路网络，可将$D_{ij}$替换为最短路距离。

在最大用车日，非零最短借还距离为 0.056 km，对应站点 S048 到 S012；非零最长借还距离为 10.100 km，对应站点 S025 到 S022。同站借还且用车超过1 min共有 325 次，这类记录可能对应临时购物、休闲骑行、试骑或绕行后回到原点，不宜简单视为异常。

\subsubsection{重点站点时刻分布}
最大日借车频次最高站点为 鹿城公共自行车站54，借车 137 次；还车频次最高站点为 鹿城公共自行车站54，还车 120 次。本文按小时统计其借还车时刻分布，并同步分析用车时长。

\fig{0.82}{../quest3/figures/q3_top_station_hourly.png}{最大日重点站点小时频次分布}
图6显示重点站点具有明显日内波动：早高峰和晚高峰是主要需求集中时段，午间和夜间存在次峰。若借车峰早于还车峰，通常对应居住区向办公区输出；若还车峰集中在晚间，则可能对应通勤回流。该分布可用于设置分时调度策略，例如早高峰前向居住区补车，晚高峰前向商业办公区补空桩。

\subsubsection{各站点高峰时段识别与归类}
将一天划分为清晨、早高峰、上午平峰、午间、下午平峰、晚高峰、夜间和深夜八个时段。站点$i$的借车高峰时段定义为
\begin{equation}
p_i^B=\arg\max_p \sum_{k:o_k=i}I(period(t_k^{borrow})=p),
\end{equation}
还车高峰时段$p_i^R$同理。进一步以站点在八个时段的借车频次和还车频次拼接为16维峰型向量，标准化后使用KMeans聚为4类。

\fig{0.78}{../quest3/figures/q3_peak_cluster_map.png}{站点峰型聚类空间分布}
图7给出了站点峰型类别的空间分布。聚类轮廓系数为 0.109。该数值不高，说明公共自行车站点峰型之间存在连续过渡，而非天然分成完全分离的类别；但四类结果仍具有运营解释价值：早高峰输出型、晚高峰回流型、全天活跃型和低强度型。对于早高峰输出型站点，应在早高峰前增加可借车辆；对于晚高峰回流型站点，应预留空桩并在峰后转运车辆。

\subsection{问题四：站点设置与锁桩配置评价}
\subsubsection{评价指标体系}
只按累计借还总量评价站点会忽略容量约束和潮汐流问题。本文设置baseline为归一化总流量：
\begin{equation}
S_i^{base}=norm(q_i),\quad q_i=\frac{B_i+R_i}{20}.
\end{equation}
主模型综合四类指标：日均借还总量$q_i$、每桩周转强度$u_i=q_i/c_i$、借还净流入率绝对值$|\rho_i|$、最大日峰值压力$h_i$。其中
\begin{equation}
\rho_i=\frac{R_i-B_i}{R_i+B_i+\epsilon}.
\end{equation}
综合重要度定义为
\begin{equation}
S_i=0.35\,norm(q_i)+0.25\,norm(u_i)+0.25\,norm(|\rho_i|)+0.15\,norm(h_i).
\end{equation}
权重设置遵循“总需求优先、容量和不平衡并重、峰值辅助”的原则。该模型输出可以直接回答站点设置和锁桩数量配置是否合理：$u_i$过高表示锁桩可能不足，$u_i$过低表示利用不足，$|\rho_i|$过高表示需要调度而不一定只靠扩容解决。

\subsubsection{评价结果}
按每桩日周转强度，本文将站点分为利用偏低、基本适宜和可能不足三类，并对借还净流入率绝对值较大的站点标记“需调度”。样本中配置评价计数为：基本适宜28个, 可能不足16个, 利用偏低16个。综合重要度前五的站点为 鹿城公共自行车站25, 鹿城公共自行车站54, 鹿城公共自行车站46, 鹿城公共自行车站30, 鹿城公共自行车站45。这些站点普遍具有较高周转强度或明显峰值压力，因此应优先扩容或强化调度。

从管理角度看，站点设置评价不能简单地将“高流量”解释为“锁桩不足”。若高流量同时伴随高周转且净流入率接近零，说明站点位置合理但容量紧张；若高流量伴随明显净流入或净流出，则说明车辆再平衡更关键；若低流量且周转低，则应检查站点可达性、周边需求或宣传引导，而不是继续扩容。

\subsection{问题五：运行规律与改进建议}
\subsubsection{运行规律总结}
基于以上结果，公共自行车服务系统表现出五条规律。第一，用车时长右偏，短途和中短途出行占主体，公共自行车主要承担接驳和短距离出行。第二，用户活跃度长尾明显，大量借车卡只在样本期使用少数几次，少数高频用户贡献较多订单。第三，站点存在潮汐流特征，居住、办公、交通枢纽等区域在早晚高峰承担不同方向的借还压力。第四，高峰压力与总流量不完全一致，部分总量中等站点因锁桩少而周转强度高。第五，同站借还不是简单噪声，可能包含休闲骑行、短途办事和临时绕行等服务场景。

\subsubsection{改进建议}
\begin{enumerate}[itemsep=0.4em]
  \item 对综合重要度高、周转强度高且净流入率不大的站点，优先增加锁桩和车辆投放，提高高峰期服务能力。
  \item 对净流入率绝对值大的站点，建立分时调度策略：早高峰前向净输出站补车，晚高峰前向净流入站预留空桩。
  \item 对利用偏低站点，结合周边道路、公交站、学校、商业设施重新评估站点可达性；若连续多期低利用，可考虑迁移或减少锁桩。
  \item 对高频用户推出通勤保障服务，例如高峰预约、常用站点余车提醒；对低频用户优化注册、押金和站点导览以提升转化率。
  \item 建议管理中心在业务系统中增加实时满桩率、空车率和调度记录字段，未来可用排队模型或仿真模型进一步评估服务失败概率。
\end{enumerate}

\section{模型检验、对比与敏感性分析}
\subsection{Baseline对比}
本文对配置评价设置了简单baseline，即只按累计流量评价站点重要性。baseline优点是直观、易解释，但缺点是不能识别容量不足和借还失衡。综合模型在总流量基础上加入周转强度、不平衡率和峰值压力，使评价结果更贴近锁桩配置和车辆调度决策。图1中，部分站点综合得分高于baseline，主要是因为其单位锁桩承载压力大；部分站点baseline较高但综合提升有限，说明其高流量可能已被较大容量消化。

\subsection{距离与约束检验}
距离计算采用经纬度球面距离，结果均非负；同站借还距离为0，在寻找最短非零距离时已剔除同站记录。用车时长清洗后限制在0至24小时内，避免负时长和跨日异常对分布产生不合理影响。配置评价中锁桩数量均为正，因此周转强度分母不为零；净流入率分母加入极小量$\epsilon$，避免零流量站点除零。

\subsection{聚类稳定性与解释性}
峰型聚类的轮廓系数为 0.109，说明类别边界并不十分清晰。这符合公共自行车站点的实际特征：站点功能通常是混合的，某些站点同时服务居住、商业和交通接驳。本文因此不把聚类结果作为硬性分类，而将其作为调度策略分组依据，并结合站点具体高峰时段表使用。

\subsection{权重敏感性分析}
为检验综合评价模型稳定性，本文分别将流量、周转、不平衡和峰值四个权重上下扰动10\%和20\%，再重新归一化权重并计算综合得分，比较Top10重点站点集合与基准模型的重合率。

\fig{0.78}{../results/sensitivity_top10_stability.png}{关键权重扰动下Top10站点稳定性}
图8显示，在所设扰动范围内Top10重点站点最低重合率为 90.0\%。这说明综合评价结论对权重变化不敏感，重点扩容和调度站点具有较强稳定性。若管理者更关注用户满意度，可适当提高峰值压力权重；若更关注资产利用率，可提高每桩周转强度权重，但核心重点站点基本保持一致。

\section{模型评价、改进与推广}
\subsection{模型优点}
第一，模型链条完整，覆盖题目要求的站点频次、用户卡、最大日、距离、峰值、配置评价和运行建议。第二，所有统计指标都可由原始记录直接复现，支撑材料中保存了代码、表格、图像和冻结数字。第三，配置评价模型不仅考虑总量，还考虑锁桩容量和借还不平衡，比单纯流量排序更贴近管理决策。第四，峰型聚类将复杂日内曲线转化为可执行调度分组，便于制定分时补车和清桩策略。

\subsection{模型局限}
最大局限是当前目录缺少真实附件，因此本文数值结果来自固定随机种子的仿真样本，不能直接作为温州市鹿城区真实运营结论。其次，距离采用站点间直线球面距离，没有考虑道路网络、骑行禁行路段和实际绕行。第三，配置评价只基于20天样本，没有纳入天气、节假日、大型活动和维修停运等外部因素。第四，锁桩配置建议使用经验阈值，没有通过实时库存仿真计算“无车可借”和“无桩可还”的概率。

\subsection{改进方向与推广}
若获得真实附件，可首先替换数据源并重算全部结果；其次引入道路网络最短路距离替代Haversine直线距离；再次结合每站实时库存和调度记录建立离散事件仿真模型，以服务失败率为目标优化锁桩和车辆调度。该模型框架也可推广到共享单车、共享电单车、校园自行车和城市微循环公交等短距离公共出行系统。

\section{结论}
本文完成了2013D公共自行车服务系统问题的完整建模流程。主要结论如下：
\begin{enumerate}[itemsep=0.35em]
  \item 站点累计借还车频次可通过日期—站点双索引稳定统计；样本中累计借车最高站点为 鹿城公共自行车站54，累计还车最高站点为 鹿城公共自行车站08。
  \item 用车时长右偏，平均 45.42 min，中位数 36.72 min，说明短途接驳是主体，但长时订单不可忽视。
  \item 用户借车卡使用次数呈长尾分布，样本中仅借一次的卡占 54.31\%，系统同时服务高频通勤用户和大量低频用户。
  \item 最大用车日为 2013-09-12，当天用车 4517 次；非零OD距离范围为 0.056 km至 10.100 km。
  \item 综合配置评价比单纯流量排序更适合指导锁桩配置和车辆调度，鹿城公共自行车站25、鹿城公共自行车站54、鹿城公共自行车站46 等站点应优先关注。
  \item 权重扰动下重点站Top10最低重合率为 90.0\%，说明评价结果具有较好稳健性。
\end{enumerate}

\section*{参考文献}
\addcontentsline{toc}{section}{参考文献}
\begin{enumerate}[label={[\arabic*]}]
\item 全国大学生数学建模竞赛组委会. 2013高教社杯全国大学生数学建模竞赛D题：公共自行车服务系统[Z]. 2013.
\item 温州市鹿城区公共自行车管理中心. 公共自行车服务系统运行资料[EB/OL]. http://www.wzbicycle.com.
\item Chrisman N. Exploring Geographic Information Systems[M]. New York: Wiley, 2002.
\item Lloyd S. Least squares quantization in PCM[J]. IEEE Transactions on Information Theory, 1982, 28(2):129--137.
\item Hyndman R J, Athanasopoulos G. Forecasting: Principles and Practice[M]. Melbourne: OTexts, 2021.
\end{enumerate}

\appendix
\section{附录：支撑材料说明}
支撑材料目录包含：\verb|code/main_modeling.py|为主程序；\verb|results/frozen_numbers.json|为论文引用的冻结数字；\verb|tables/|保存各问题结果表；\verb|quest1|、\verb|quest2|、\verb|quest3|分别保存分问题输出和图表；\verb|data/|保存题面文档副本与本次分析使用的数据表。运行方式为：
\begin{verbatim}
python 支撑材料/code/main_modeling.py
\end{verbatim}
编译论文方式为：
\begin{verbatim}
cd 支撑材料/papper
xelatex -interaction=nonstopmode 论文.tex
xelatex -interaction=nonstopmode 论文.tex
xelatex -interaction=nonstopmode 论文.tex
\end{verbatim}

\section{附录：三层质量门控状态}
\begin{itemize}
\item L1建模合理性：通过。每个子问题均给出题目要求、模型输出和解释；配置评价设置了baseline与综合模型对比。
\item L2求解正确性：有条件通过。代码可复现、随机种子固定、结果冻结；但由于真实附件缺失，数值为仿真样本结果，不能替代真实运营结论。
\item L3论文质量：通过。论文包含摘要数字、公式、图表解释、模型检验、敏感性分析、评价与支撑材料说明。
\end{itemize}
\end{document}
